\section{The Logic of Paradox}
\label{sec:many-valued-paraconsistent-04}

Use the same value set and operations, but designate
\[
  D_{\mathsf{LP}}=\set{\tfrac12,1}.
\]
Interpret \(\tfrac12\) as carrying both positive and negative support.  The
resulting matrix defines the Logic of Paradox \(\mathsf{LP}\).  Its
value-labelled tableau is sound and complete by the finite-matrix theorem.

\begin{definition}[Paraconsistent consequence relation]
A consequence relation is paraconsistent iff there are formulas
\(A,B\) such that
\[
  A,\neg A\notentails B.
\]
This definition denies unrestricted explosion.  It does not assert that any
particular contradiction is derivable or true.
\end{definition}

\begin{example}[Counterexample valuation to explosion in \(\mathsf{LP}\)]
Set
\[
  v(p)=\tfrac12,
  \qquad
  v(q)=0.
\]
Both \(p\) and \(\neg p\) receive the designated value \(\tfrac12\), while
\(q\) is nondesignated.  Hence
\[
  p,\neg p\notentails_{\mathsf{LP}}q.
\]
At the same valuation, \(p\lor\neg p\) is designated.  Paraconsistency and the
failure of excluded middle are therefore independent properties; compare the
two choices of designated values on the same algebra.
\end{example}

\begin{warning}[Truth values are formal data]
Labels such as ``both'' or ``indeterminate'' are interpretations of matrix
values.  The matrix consequence relation is determined by the truth-value set,
operations, and designated set.  An informal label cannot replace those data.
\end{warning}

